\scanpage{37}
\mspage{- \handwritten{31} -}
% archive sheet number: 36

\begin{center}
CHAPITRE II (Etat 1)

FONCTIONS HOLOMORPHES D'UNE VARIABLE COMPLEXE.

\handwritten{\S}1. Le principe du prolongement analytique.
\end{center}

\uline{\handwritten{1}}. \uline{Prolongement analytique \unsure{.|}d'une fonction holomorphe}.

Dans ce \handwritten{\S}\unsure{,|;} K désignera \uline{le corps des nombres réels ou le corps des nombres complexes}, E un espace normé complet sur K.

PROPOSITION 1.- \uline{Dans l'espace} $K^p$, \uline{soient} P \uline{et} Q \uline{deux polycylindres ouverts}, \uline{de centres} $\handwritten{a}$ \uline{et} $\handwritten{b}$, \uline{tels que} $P\handwritten{\cap}Q$ \uline{ne soit pas vide}. \uline{Soient} $\handwritten{f}(\handwritten{x})=\handwritten{\sum c}_{n_1n_2..n_p}(x_1-a_1)^{n_1}\ldots(x_p-a_p)^{n_p}$ \uline{la somme d'une série entière par rapport aux} $x_i-a_i$, \uline{sommable dans} P, \uline{et} $\handwritten{g(}\deleted{\handwritten{g}}\inserted{\handwritten{x}}\handwritten{)}\handwritten{=}$ $=\handwritten{\sum d}_{n_1n_2..n_p}(x_1-b_1)^{n_1}\ldots(x_p-b_p)^{n_p}$ \uline{la somme d'une série entière par rapport aux} $x_i-b_i$, \uline{sommable dans} Q. \uline{S'il existe un ensemble ouvert non vide} $U\handwritten{\subset}P\handwritten{\cap}Q$ \uline{tel que} $\handwritten{f}(\handwritten{x})=\handwritten{g}(\handwritten{x})$ \uline{pour tout} $\handwritten{x\in}U$, \uline{on a} $\handwritten{f}(\handwritten{x})=\handwritten{g}(\handwritten{x})$ \uline{pour} tout $\handwritten{x\in}P\handwritten{\cap}Q$.

Soit en effet $\handwritten{x}$ un point de U et $\handwritten{y}$ un point quelconque de $P\handwritten{\cap}Q$ ; le segment fermé d'extrémités $\handwritten{x}$ et $\handwritten{y}$ appartient tout entier à $P\handwritten{\cap}Q$, qui est convexe ; la fonction $\handwritten{h}(t)=\handwritten{f}(\handwritten{x}+t(\handwritten{y}-\handwritten{x}))-\handwritten{g}(\handwritten{x}+t(\handwritten{y}-\handwritten{x}))$ est donc une fonction holomorphe de la variable réelle t dans un intervalle ouvert I contenant $\handwritten{[}0,1\handwritten{]}$. Désignons par $\handwritten{\rho}$ la borne supérieure de l'ensemble A des t tels que $0\handwritten{\leq}t\handwritten{\leq}1$ et que $\handwritten{h}(s)=0$ pour $0\handwritten{\leq}s\handwritten{\leq}t$ ; on a par hypothèse $\handwritten{\rho>}0$, et il est clair par continuité que $\handwritten{\rho\in}A$ ; nous allons montrer que $\handwritten{\rho}=1$ ce qui démontrera la proposition. Dans le cas contraire, la fonction $\handwritten{h}(t)$ serait égale à la somme d'une série entière en $t-\handwritten{\rho}$, convergente dans un intervalle ouvert de centre $\handwritten{\rho}$ ; mais on a $\handwritten{h}(t)=0$ pour tout t tel que $0\handwritten{\leq}t\handwritten{\leq\rho}$, donc d'après le principe des zéros isolés, les coefficients de cette série en $t-\handwritten{\rho}$ sont tous nuls on a donc aussi $\handwritten{h}(t)=0$ pour $\handwritten{\rho<}t\handwritten{<\rho}+\handwritten{\alpha}$, où $\handwritten{\alpha>}0$ ; mais cela est
